\section{对角化失败与 Jordan 形：独立特征方向不够时怎么办}
\label{sec:03-Linear-Algebra/05-Eigenvalues-and-Eigenvectors/03}

一个最简单的积分器：状态由位置和速度组成，每一步速度不变，位置加上速度，

\[
p_{k+1}=p_k+v_k,\qquad v_{k+1}=v_k,
\qquad
A=\begin{pmatrix}1&1\\0&1\end{pmatrix}.
\]

按上一节的判据检查稳定性。特征多项式是 \((1-\lambda)^2\)，唯一特征值 \(\lambda=1\)，谱半径恰好等于 \(1\)，看起来不增不减。但只要 \(v_0\neq 0\)，位置就按 \(p_k=p_0+kv_0\) 一路漂走，任何有界的结论都是错的。

问题不在特征值。问题在于上一节的结论有个前提——把 \(x_0\) 拆成特征方向的组合——而这个矩阵拆不了。\(A-I=\begin{pmatrix}0&1\\0&0\end{pmatrix}\) 的零空间只有 \((1,0)^{\trans}\) 这一条线：代数重数说好两个，几何重数只给一个。缺的那一维不会凭空消失，它换了个形式重新出现在解里。

\subsection{只有一条不变直线，其余方向被推着走}

先把这个矩阵的作用看清楚。它保持每个点的纵坐标，按纵坐标的大小把横坐标推一段——这是剪切。

\begin{center}
\begin{tikzpicture}[x=1cm,y=1cm,>=stealth,font=\small]
  \begin{scope}
    \draw[gray!55,->] (-1.9,0) -- (2.4,0);
    \draw[gray!55,->] (0,-0.5) -- (0,2.0);
    \draw[black!30,dashed] (-1.9,0) -- (2.4,0);
    \draw[->,black!45,thick] (0,0) -- (0,1.4);
    \draw[->,blue!70!black,thick] (0,0) -- (1.4,1.4);
    \draw[->,black!45,thick] (0,0) -- (-1.2,0.9);
    \draw[->,blue!70!black,thick] (0,0) -- (-0.3,0.9);
    \draw[->,red!65!black,very thick] (0,0) -- (1.7,0);
    \node[below,red!65!black] at (1.7,-0.12) {不变方向};
    \node[above left,black!55] at (0,1.4) {\(u\)};
    \node[above right,blue!70!black] at (1.4,1.4) {\(Au\)};
    \node at (0.3,-1.0) {剪切：横坐标按纵坐标推移};
  \end{scope}
  \begin{scope}[shift={(5.6,0)}]
    \draw[gray!55,->] (-0.4,0) -- (4.0,0);
    \draw[gray!55,->] (0,-0.5) -- (0,2.0);
    \draw[black!30,dashed] (-0.4,1) -- (4.0,1);
    \foreach \k in {0,1,2,3,4}
      \fill[blue!70!black] (0.85*\k,1) circle (2pt);
    \draw[->,blue!45!black] (0.1,1.25) -- (0.75,1.25);
    \draw[->,blue!45!black] (0.95,1.25) -- (1.6,1.25);
    \draw[->,blue!45!black] (1.8,1.25) -- (2.45,1.25);
    \node[below] at (0,0.85) {\(k=0\)};
    \node[below] at (3.4,0.85) {\(k=4\)};
    \node at (1.8,-1.0) {每步向右移动一个固定量};
  \end{scope}
\end{tikzpicture}

\emph{左：只有横轴上的向量方向不变，其余向量都被推离原来的直线。右：起点 \((0,1)^{\trans}\) 反复作用后沿横轴匀速前进——特征值是 \(1\)，位移却线性累积。}
\end{center}

代数上把这件事说穿只要一行。写 \(A=I+N\)，其中

\[
N=\begin{pmatrix}0&1\\0&0\end{pmatrix},\qquad N^2=0 .
\]

\(I\) 与 \(N\) 交换，二项式展开合法，而 \(j\ge 2\) 的项全为零，于是

\[
A^k=(I+N)^k=I+kN=\begin{pmatrix}1&k\\0&1\end{pmatrix}.
\]

指数因子 \(1^k\) 什么也没做，多出来的 \(k\) 才是全部内容。它来自图中那个反复推移的动作：\((0,1)^{\trans}\) 不是特征向量，每作用一次就在横坐标上加一个单位，累积 \(k\) 次就是 \(k\)。用状态的语言说，速度被持续积分进位置，一个状态量在不断吃进另一个状态量。

值得强调的是这类矩阵有多常见：位置与速度的更新、离散积分器、带反馈的低通滤波、临界阻尼的振子，都是这个形状。对角化失败不是需要绕开的病例，而是一类必须能读懂的结构。

\subsection{广义特征向量把缺掉的方向补回来}

真正的特征向量满足 \((A-\lambda I)v_1=0\)，被 \(A-\lambda I\) 一次消灭。方向不够时退而求其次：找一个向量，被 \(A-\lambda I\) 作用后不归零，但恰好落在已有的特征方向上，

\[
(A-\lambda I)v_2=v_1 .
\]

再作用一次就归零了，因为 \((A-\lambda I)^2v_2=(A-\lambda I)v_1=0\)。这样的 \(v_2\) 叫广义特征向量。上面的例子里取 \(v_1=(1,0)^{\trans}\)、\(v_2=(0,1)^{\trans}\) 即可，而这正是图中那个被反复推移的向量：它自己保不住方向，却把每一步的推移量交给了 \(v_1\)。

更长的链遵循同样的递推 \((A-\lambda I)v_{j+1}=v_j\)，链首是真正的特征向量，链上每一环被 \(A-\lambda I\) 推向前一环，直到链首被消灭。Jordan 标准形定理断言，任何复方阵都相似于一个由这种链拼成的分块对角矩阵，块的形状是

\[
J_r(\lambda)=
\begin{pmatrix}
\lambda&1&&\\
&\lambda&\ddots&\\
&&\ddots&1\\
&&&\lambda
\end{pmatrix}
\quad(r\times r),
\qquad
A=SJS^{-1}.
\]

对角线上的 \(\lambda\) 是缩放，超对角线上的每个 \(1\) 是一条链节。手工求出 \(S\) 的过程是机械的线性方程求解，这里略过，因为它既不影响判断也不是实际算法（本节末尾会说明为什么）。要记住的只有一条对应：\(A\) 可对角化，当且仅当所有 Jordan 块的大小都是 \(1\)，也就是超对角线上没有 \(1\)。

\subsection{链节产生的是多项式因子}

把块写成 \(J=\lambda I+N\)，\(N\) 只在超对角线上有 \(1\)，是幂零的：块大小为 \(r\) 时 \(N^r=0\)。\(\lambda I\) 与一切交换，所以

\[
J^k=(\lambda I+N)^k=\sum_{j=0}^{r-1}\binom{k}{j}\lambda^{k-j}N^j,
\]

级数在 \(j=r-1\) 处自动截断。除了指数因子 \(\lambda^k\)，还冒出 \(\binom{k}{1}=k\)、\(\binom{k}{2}\approx k^2/2\) 这些多项式因子，最高次数是 \(r-1\)。连续时间是同一回事：

\[
e^{Jt}=e^{\lambda t}\Bigl(I+tN+\frac{t^2}{2!}N^2+\cdots+\frac{t^{r-1}}{(r-1)!}N^{r-1}\Bigr).
\]

微分方程课上出现的 \(te^{\lambda t}\)、\(t^2e^{\lambda t}\) 就是从这里来的。重根时``解要乘上 \(t\)''不是凑出来的特解形式，而是 Jordan 块强制的结果——一个大小为 \(r\) 的块，意味着某个状态量被反复积分了 \(r-1\) 次。

于是可以给出通项的形状：可对角化时轨道由 \(\lambda_i^k\) 或 \(e^{\lambda_it}\) 组成，一般情形则是 \(k^{j}\lambda_i^k\) 或 \(t^je^{\lambda_it}\)，\(j\) 不超过对应块大小减一。

\subsection{临界情形因此没有安全余量}

有了多项式因子，上一节留下的悬案就能结清。离散系统在 \(\rho(A)<1\) 时衰减，\(\rho(A)>1\) 时发散，这两条与 Jordan 结构无关——多项式再怎么长，也压不过一个模小于 \(1\) 的数的幂，也救不了一个模大于 \(1\) 的数的幂。真正需要额外信息的只有 \(\rho(A)=1\)：模等于 \(1\) 的特征值若都对应大小为 \(1\) 的块，系统保持有界；只要有一个块大于 \(1\)，\(k^{j}\) 就无人制约，系统按多项式速度发散。本节开头那个积分器正是最小的例子。连续时间同理，临界的是 \(\operatorname{Re}\lambda=0\) 的特征值。

这条边界在实践中很硌人。想让循环网络既不遗忘也不爆炸，自然的想法是把谱半径调到 \(1\) 附近；可 \(\rho=1\) 恰好落在多项式增长可能出现的位置上，缺陷结构会让隐状态随步数缓慢累积，而不是稳稳地保持。门控机制的价值也在这里：它不是把某个特征值精确钉在 \(1\)，而是让网络按内容决定每一步保留多少，把是否累积交给数据而不是交给谱的临界值。

同样的谨慎适用于 Markov 链。模为 \(1\) 的特征值不止 \(\lambda=1\) 一个时，链可能是周期的，分布长期振荡而不收敛；\hyperref[sec:06-Random-Processes/03-Markov-Chains/02]{``可达、通信类与周期性''}讨论的非周期性条件，落到谱上就是要求单位圆上除 \(1\) 之外没有别的特征值。

\subsection{最小多项式看的是最大的那个块}

特征多项式记录了每个特征值的代数重数，但它分不清``三个大小为 \(1\) 的块''和``一个大小为 \(3\) 的块''——两种情形的特征多项式都是 \((t-\lambda)^3\)，动力学却完全不同。

最小多项式 \(m_A\) 补上了这个区分。它是使 \(m_A(A)=0\) 的最低次首一多项式，而对每个特征值，因子 \((t-\lambda)\) 的指数恰好等于该特征值对应的最大块的大小。理由不难看：在单块 \(J_r(\lambda)\) 上 \((J_r(\lambda)-\lambda I)^s=N^s\)，它为零当且仅当 \(s\ge r\)；要让所有块同时归零，指数必须覆盖最大的 \(r\)，取最低次就正好等于它。由此得到一条比数特征值更准的判据：

\[
A\ \text{可对角化}
\iff
m_A(t)\ \text{无重根}.
\]

特征值尽管重复，只要每个特征值的块都是 \(1\times 1\)，最小多项式就没有重根，对角化照样成立。

\subsection{理论上的标准形，数值上的雷区}

Jordan 形在理论上是相似分类的完整答案，在浮点运算里却几乎不能用，因为它对扰动极不稳定。看

\[
A_\varepsilon=\begin{pmatrix}1&1\\ \varepsilon&1\end{pmatrix}.
\]

\(\varepsilon=0\) 时它是一个 \(2\times 2\) 的 Jordan 块；只要 \(\varepsilon>0\)，特征值立刻分裂成 \(1\pm\sqrt{\varepsilon}\)，矩阵变得可以对角化。注意分裂的幅度是 \(\sqrt{\varepsilon}\) 而不是 \(\varepsilon\)：双精度下一个 \(10^{-16}\) 量级的舍入误差，会把重特征值撑开约 \(10^{-8}\)，比机器精度大了八个数量级。同时两个特征向量从几乎平行变成严格无关，\(S\) 病态得离谱。

结果是，你无法在浮点数里可靠地判断一个特征值究竟是严格重根还是两个挨得很近的根，而 Jordan 形的全部结构恰恰押在这个判断上。数值软件因此改用 Schur 分解：用酉相似把矩阵化成上三角，对角线上仍是特征值，变换阵正交因而稳定；需要块结构时再在此基础上做谨慎的聚类。\hyperref[sec:08-Numerical-Analysis/05-Numerical-Linear-Algebra/05]{``从幂迭代到 QR 特征值算法''}讲的就是这条路线。

这里浮现出一条更一般的分工：数学上最简的标准形，与数值上最可靠的表示，常常不是同一个对象。前者回答``若精确成立会怎样''，后者回答``在舍入误差下还能保证什么''。同一句话在\hyperref[sec:03-Linear-Algebra/07-SVD-and-Data-Applications/02]{``伪逆与病态问题''}里还会以秩的形式再出现一次。

\subsection{与 SVD 是两个不同的问题}

既然特征分解会因方向不足而失败、会因基几乎平行而病态、还会跑到复数域，为什么不干脆用 SVD？因为两者回答的问题不同。

特征分解问的是同一个空间内部的不变方向：输入与输出住在一起，哪些方向只被缩放。这个问题本身就可能无解，Jordan 块就是无解时的记录。SVD 允许输入与输出各用一组自己的正交基，不要求方向留在原地，只要求两侧各自正交、中间是非负伸缩，于是对任意矩阵——包括长方形的——总是存在。

所以 Jordan 形回答``一个方阵在相似变换下最简的样子是什么''，SVD 回答``一个线性映射把哪些正交输入方向送到哪些正交输出方向''。研究反复迭代、长期演化、稳定性，问的是前一个；研究一次映射的放大倍数、低秩结构、最小二乘，问的是后一个。若发现两个都想要，多半是两个问题被混在了一起，值得先分开再动手。

眼下还有一件事没做完。本节和上一节分别算出了 \(A^k\) 与 \(e^{At}\)，方法看着像是两次巧合：一次是对角矩阵逐项取幂，一次是幂级数逐项代入。下一节把它们放进同一个框架，问一个更大的问题——把任意标量函数作用到矩阵上，究竟该怎么定义。
